% !TeX root = ../main.tex
%
% Necesită în preambul:
%   \usepackage{booktabs, graphicx, amsmath, amssymb}
%   \usepackage[ruled,vlined]{algorithm2e}   % sau algorithm + algpseudocode
%   \usepackage{listings}
%   \usepackage[utf8]{inputenc}              % pentru diacriticele româneşti (dacă se compilează cu pdfLaTeX)
%   \usepackage[romanian]{babel}
%
\chapter{Abordarea investigată}
\label{chapter:proposedApproach}

\section{Prezentare generală}
\label{sec:overview}

Acest capitol documentează întreaga succesiune de experimente pe care le-am
realizat asupra MLEP (Multi-granularity Local Entropy Patterns)~\cite{yuan2025mlep},
un detector de imagini generate de inteligenţa artificială al cărui etaj de
intrare înlocuieşte imaginea brută cu un teanc de hărţi de entropie locală.
Obiectivul nostru a fost dublu: în primul rând, reproducerea detectorului
publicat şi confirmarea comportamentului raportat pe date curate; în al doilea
rând, găsirea unei modificări a etajului său de entropie care să îmbunătăţească
fie acurateţea brută de detecţie, fie --- mai important --- robusteţea la
degradările pe care orice imagine le întâlneşte în lumea reală: neclaritatea
(blur) şi recompresia cu pierderi.

Raportăm fiecare experiment, inclusiv pe cele care au eşuat, deoarece tiparul
eşecurilor este el însuşi rezultatul. Capitolul este organizat ca un lanţ în
care fiecare experiment este motivat de rezultatul celui precedent:

\begin{enumerate}
  \item Reproducem modelul de referinţă şi, făcând acest lucru, descoperim un
        defect al \emph{statisticilor BatchNorm} care invalidează evaluarea
        naivă (Secţiunea~\ref{sec:setup-bn}).
  \item Variem \emph{dimensiunea ferestrei glisante} a operatorului de entropie,
        cel mai evident hiperparametru (Secţiunea~\ref{sec:exp-window}).
        Ferestrele mai mari dăunează monoton.
  \item Deoarece nicio fereastră individuală nu a ajutat, testăm
        \emph{combinarea} mai multor ferestre, precum şi două corecţii aduse
        acestei combinaţii --- normalizarea per fereastră şi înregistrarea
        spaţială exactă (Secţiunea~\ref{sec:exp-multiwindow}). Niciuna nu ajută.
  \item Deoarece \emph{geometria} ferestrei a fost epuizată, ne întrebăm dacă
        \emph{funcţionala} de entropie însăşi este alegerea greşită şi testăm
        entropiile Rényi, Tsallis şi de permutare, atât într-un etaj de
        trăsături clasice, cât şi într-un etaj profund
        (Secţiunea~\ref{sec:exp-entropy}). Toate se încadrează în zgomot.
  \item Deoarece reponderările scalare ale aceleiaşi distribuţii de valori nu au
        schimbat nimic, testăm o statistică pe care marginalele, demonstrabil,
        nu o pot exprima --- \emph{entropia comună de culoare între canale}
        (Secţiunea~\ref{sec:exp-color}). Niciun câştig.
  \item Deoarece toate variantele etajului de intrare au eşuat pe date
        \emph{curate}, atacăm robusteţea în mod direct: o \emph{separare a
        texturilor bogate/sărace} în stil PatchCraft
        (Secţiunea~\ref{sec:exp-texture}) şi \emph{antrenarea cu augmentare prin
        degradări} (Secţiunea~\ref{sec:exp-aug}). Separarea texturilor nu ajută;
        augmentarea ajută doar în cadrul degradării exacte pe care a fost
        antrenată.
  \item În final ne întrebăm dacă etajul de intrare poate măcar \emph{raporta}
        degradarea pe care o suferă, printr-un model cu mai multe capete
        antrenat 200k paşi (Secţiunea~\ref{sec:exp-degradation}). Acesta poate
        identifica degradarea aproape perfect, rămânând totuşi incapabil să
        detecteze prin ea.
\end{enumerate}

Secţiunea~\ref{sec:conclusion} formulează concluzia care motivează restul
lucrării: nicio intervenţie asupra etajului de entropie nu a recuperat
robusteţea, aşa că ne îndreptăm către combinarea MLEP cu un al doilea detector,
diferit din punct de vedere arhitectural.

\section{Algoritmul de referinţă}
\label{sec:algorithm}

\subsection{Intuiţia}

MLEP se bazează pe o observaţie simplă privind modul în care modelele generative
sintetizează pixelii. Un senzor de cameră produce zgomot de fotoni şi de citire
independent în fiecare punct, astfel încât pixelii vecini ai unei fotografii
naturale aproape niciodată nu poartă valori pe 8 biţi exact egale. Un generator,
prin contrast, îşi produce ieşirea prin convoluţii transpuse sau printr-un
decodor care supraeşantionează un cod latent de rezoluţie mică, iar
supraeşantionarea \emph{duplică} valori: o expandare de tip nearest-neighbour
sau biliniară face ca pixelii adiacenţi să fie egali sau cvasi-egali mult mai
des decât ar face-o un senzor.

MLEP transformă acest lucru într-o trăsătură renunţând complet la conţinutul
fotometric şi păstrând doar o măsură a \emph{diversităţii locale de valori}:
pentru fiecare poziţie din imagine numără câte valori distincte apar într-o
fereastră glisantă mică. Reţeaua nu vede niciodată luminozitatea sau culoarea
--- doar tiparul spaţial al diversităţii locale, ,,tiparul de entropie locală''.

\subsection{Entropia locală}

Fie $x \in \mathbb{R}^{C \times H \times W}$ imaginea (normalizată ImageNet) şi
fie $\mathcal{W}_{ij}^{(w)}$ fereastra $w \times w$ al cărei colţ stânga-sus este
$(i,j)$, conţinând $K = w^2$ valori. Notând cu $p_g$ fracţiunea celor $K$ valori
egale cu valoarea distinctă $g$, entropia Shannon a ferestrei este
\begin{equation}
  H_{\text{S}}(\mathcal{W}) = -\sum_{g} p_g \log_2 p_g .
  \label{eq:shannon}
\end{equation}
Harta este calculată cu pasul $1$ şi \emph{independent pe fiecare canal de
culoare}, astfel încât un operator $w\times w$ aplicat unei imagini
$C\times H\times W$ produce $C \times (H-w+1) \times (W-w+1)$ valori.

Pentru valoarea implicită $w=2$ fereastra conţine doar $K=4$ pixeli, deci
distribuţia de valori poate lua exact cinci forme, iar MLEP exploatează acest
lucru printr-o tabelă de căutare, nu printr-un calcul general.
Tabelul~\ref{tab:2x2} le enumeră. De remarcat că implementarea publicată
stochează $0.8$ pentru cazul ,,trei egale'', o rotunjire a valorii exacte
$\log_2 4 - \tfrac{3}{4}\log_2 3 = 0.8113$; păstrăm această rotunjire pentru ca
reproducerea noastră să fie compatibilă la nivel de bit cu ponderile publicate.

\begin{table}[t]
  \centering
  \caption{Cele cinci configuraţii ale unei ferestre $2\times2$ şi entropia pe
  care MLEP le-o atribuie. ,,Perechi egale'' numără, dintre cele şase perechi
  neordonate ale celor patru pixeli, câte sunt egale --- acest număr întreg
  identifică unic cazul şi este ceea ce face posibilă tabela de căutare cu cinci
  intrări.}
  \label{tab:2x2}
  \begin{tabular}{llcc}
    \toprule
    Caz & Dimensiunile grupurilor & Perechi egale & $H_{\text{S}}$ \\
    \midrule
    toate patru egale              & $(4)$       & 6 & 0.0 \\
    trei egale, una diferită       & $(3,1)$     & 3 & 0.8 \\
    două perechi                   & $(2,2)$     & 2 & 1.0 \\
    o pereche, două valori unice   & $(2,1,1)$   & 1 & 1.5 \\
    toate patru diferite           & $(1,1,1,1)$ & 0 & 2.0 \\
    \bottomrule
  \end{tabular}
\end{table}

\subsection{Multi-granularitate: piramida de scări}

O singură dimensiune de fereastră vede structura la o singură frecvenţă
spaţială. Prin urmare, MLEP calculează harta de entropie la trei \emph{scări}.
Scara $s$ este realizată prin subeşantionarea biliniară a imaginii cu factorul
$s$ şi supraeşantionarea imediată înapoi la dimensiunea originală, ceea ce
distruge detaliile mai fine de $1/s$ pixeli, păstrând totodată forma tensorului
fixă:
\begin{equation}
  x^{(s)} = \text{Up}_{H\times W}\big(\text{Down}_{s}(x)\big),
  \qquad s \in \{1.0,\, 0.5,\, 0.25\}.
\end{equation}
Cu $C=3$ canale, $|S| = 3$ scări şi o singură dimensiune de fereastră, etajul de
intrare emite $3 \times 3 \times 1 = 9$ canale, exact ceea ce consumă prima
convoluţie a reţelei. În general
\begin{equation}
  C_{\text{in}} = |\text{pânze}| \cdot |S| \cdot |\mathcal{W}| \cdot |\mathcal{M}| \cdot 3,
  \label{eq:channels}
\end{equation}
unde $\mathcal{W}$ este mulţimea dimensiunilor de fereastră, iar $\mathcal{M}$
mulţimea funcţionalelor de entropie --- ambele fiind variate în experimentele de
mai jos.

\subsection{Amestecarea blocurilor}

Înainte de calculul hărţilor de entropie, imaginea este tăiată în blocuri
$2\times2$ neîntrepătrunse, care sunt apoi permutate aleator pe întregul cadru.
Acest lucru distruge semantica globală --- forma obiectelor, dispunerea,
conţinutul scenei --- lăsând însă intacte relaţiile dintre pixelii din interiorul
unui bloc, ceea ce forţează clasificatorul să se bazeze pe statistici locale şi
nu pe ceea ce înfăţişează imaginea. În timpul antrenării se extrage o permutare
nouă la fiecare trecere înainte; la evaluare permutarea este fixată de un
generator cu sămânţă, astfel încât scorarea să fie deterministă.

\begin{quote}\small
\textbf{Notă de implementare.} Codul publicat apela
\texttt{torch.manual\_seed(99)} în interiorul trecerii înainte. Acest lucru atât
îngheţa amestecarea la un singur tipar în timpul antrenării, cât şi reiniţializa
generatorul de numere aleatoare \emph{global} la fiecare pas, distrugând
aleatorismul folosit pentru formarea loturilor şi pentru augmentare. L-am
înlocuit cu un \texttt{torch.Generator} local, folosit doar la evaluare. Aceasta
este o corecţie a infrastructurii experimentale, nu o modificare a metodei, şi
este motivul pentru care curbele noastre de antrenare diferă de cele ale unei
rulări naive a codului publicat.
\end{quote}

\subsection{Reţeaua de bază}

Teancul de entropie este dat unui ResNet trunchiat: \texttt{conv1} este
redimensionat pentru a accepta $C_{\text{in}}$ canale în loc de $3$, iar dintre
straturi se păstrează doar \texttt{layer1} şi \texttt{layer2} înainte de
agregarea globală prin medie şi de un singur cap liniar. Rezultatul este un
model mic --- $0.68$M parametri pentru varianta ResNet-18 folosită în baleiaje
şi $1.44$M pentru varianta ResNet-50 folosită la rularea lungă finală ---
deoarece munca discriminativă este făcută de etajul de intrare, nu de adâncime.

\subsection{Pseudocod}

\begin{algorithm}[t]
\caption{Trecerea înainte MLEP (generalizată; modelul publicat este cazul
$\mathcal{W}=\{2\}$, $\mathcal{M}=\{\text{shannon}\}$, fără separare de texturi)}
\label{alg:mlep}
\KwIn{imaginea $x \in \mathbb{R}^{3\times H\times W}$; scările $S$; ferestrele
$\mathcal{W}$; modurile de entropie $\mathcal{M}$}
\KwOut{logit(-ii) $z$}
\BlankLine
$\mathcal{X} \leftarrow \textsc{TextureSplit}(x)$ \tcp*{$\{x\}$, sau (bogat, sărac); Sec.~\ref{sec:exp-texture}}
\ForEach{pânza $c \in \mathcal{X}$}{
  $c \leftarrow \textsc{ShuffleBlocks}(c, 2)$ \tcp*{aleator la antrenare, cu sămânţă la evaluare}
}
$F \leftarrow [\;]$\;
\ForEach{pânza $c \in \mathcal{X}$}{
  \ForEach{scara $s \in S$}{
    $c_s \leftarrow \textsc{Up}_{H\times W}(\textsc{Down}_s(c))$ \tcp*{identitate dacă $s=1$}
    \ForEach{fereastra $w \in \mathcal{W}$}{
      \ForEach{modul $m \in \mathcal{M}$}{
        $P \leftarrow \textsc{Unfold}(c_s, w, \text{pas}=1)$ \tcp*{$(3, (H\!-\!w\!+\!1)(W\!-\!w\!+\!1), w^2)$}
        $E \leftarrow \textsc{Entropy}_m(P)$ \tcp*{pe canal, pe fereastră}
        \lIf{$w=2$ \textbf{şi} $m=$ shannon}{$E \leftarrow \textsc{Lookup}_5(P)$ \tcp*{Tabelul~\ref{tab:2x2}}}
        \lIf{normalizare}{$E \leftarrow E / H_{\max}(m,w)$}
        \lIf{aliniere $=$ pad}{$E \leftarrow \textsc{ReplicatePad}(E, (w-w_{\min})/2)$}
        adaugă $E$ la $F$\;
      }
    }
  }
}
$F \leftarrow \textsc{ResizeToCommon}(F)$;\quad $z_0 \leftarrow \textsc{Concat}(F)$ \tcp*{$C_{\text{in}}$ canale, Ec.~\ref{eq:channels}}
$z \leftarrow \textsc{Head}(\textsc{Layer2}(\textsc{Layer1}(\textsc{Stem}(z_0))))$\;
\Return $z$\;
\end{algorithm}

\subsection{Un exemplu lucrat}
\label{sec:worked-example}

Pentru a face concret mecanismul --- şi modul său de eşec --- să considerăm o
regiune $4\times4$ dintr-un singur canal de culoare. Luăm mai întâi o porţiune
de cer fotografiat real, unde zgomotul senzorului perturbă fiecare punct:

\[
X_{\text{real}} =
\begin{pmatrix}
120 & 121 & 119 & 122\\
121 & 123 & 120 & 118\\
119 & 120 & 122 & 121\\
122 & 118 & 121 & 123
\end{pmatrix}
\]

Glisarea unei ferestre $2\times2$ cu pasul 1 produce o hartă de entropie
$3\times3$. Fereastra de la poziţia $(0,0)$ conţine $\{120,121,121,123\}$: o
pereche şi două valori unice, deci $1.5$ conform Tabelului~\ref{tab:2x2}.
Fereastra de la $(0,1)$ conţine $\{121,119,123,120\}$, toate patru diferite,
deci $2.0$. Continuând:

\[
E_{\text{real}} =
\begin{pmatrix}
1.5 & 2.0 & 2.0\\
2.0 & 1.5 & 2.0\\
2.0 & 2.0 & 1.5
\end{pmatrix},
\qquad \overline{E}_{\text{real}} = 1.833 .
\]

Luăm acum regiunea corespunzătoare dintr-o imagine GAN, unde un pas de
supraeşantionare $2\times$ a duplicat fiecare valoare într-o dală $2\times2$:

\[
X_{\text{fake}} =
\begin{pmatrix}
120 & 120 & 122 & 122\\
120 & 120 & 122 & 122\\
124 & 124 & 126 & 126\\
124 & 124 & 126 & 126
\end{pmatrix}
\Longrightarrow
E_{\text{fake}} =
\begin{pmatrix}
0.0 & 1.0 & 0.0\\
1.0 & 2.0 & 1.0\\
0.0 & 1.0 & 0.0
\end{pmatrix},
\qquad \overline{E}_{\text{fake}} = 0.667 .
\]

Separarea este mare şi, esenţial, \emph{structurată}: harta falsului nu este
doar mai mică în medie, ci este periodică, cu zerouri oriunde o fereastră cade
în interiorul unei dale supraeşantionate şi cu maxime oriunde traversează colţul
unei dale. Această reţea periodică este ceea ce învaţă reţeaua convoluţională de
bază şi este motivul pentru care MLEP atinge o acurateţe aproape perfectă pe
imagini GAN curate.

\paragraph{De ce se prăbuşeşte acest lucru la recompresie.}
Aplicăm acum o compresie JPEG puternică asupra porţiunii \emph{reale}. Pasul de
cuantizare DCT aplatizează fluctuaţiile mici, iar după decuantizare un bloc
$8\times8$ de pixeli de cer aproape egali returnează de regulă o mână de valori
repetate:

\[
X_{\text{real}}^{\text{JPEG}} =
\begin{pmatrix}
120 & 120 & 120 & 120\\
120 & 120 & 120 & 120\\
122 & 122 & 122 & 122\\
122 & 122 & 122 & 122
\end{pmatrix}
\Longrightarrow
\overline{E}_{\text{real}}^{\text{JPEG}} = 0.333 .
\]

Porţiunea reală are acum un scor \emph{mai mic} decât porţiunea falsă neatinsă
($0.333 < 0.667$). Trăsătura nu a devenit doar zgomotoasă --- ordonarea ei
s-a \textbf{inversat}. Un clasificator antrenat să citească ,,entropie locală
mică $\Rightarrow$ generată'' va eticheta acum, cu încredere, fotografii reale
comprimate drept generate.

Această singură observaţie prezice, cantitativ, rezultatul dominant al fiecărui
experiment din acest capitol: precizia medie în condiţie coruptă a fiecărei
variante MLEP pe care am antrenat-o se situează \emph{sub} $0.5$, în intervalul
$0.41$--$0.49$. O precizie medie sub nivelul hazardului nu înseamnă performanţă
slabă; înseamnă o ordonare inversată. Etajul de entropie nu poate distinge
,,neted pentru că scena este netedă'' de ,,neted pentru că un codec l-a
netezit'', iar fiecare intervenţie descrisă mai jos este, într-un fel sau altul,
o încercare de a-i oferi un canal în care această distincţie să poată fi
exprimată.

\section{Configuraţia experimentală}
\label{sec:setup}

\subsection{Hardware şi software}

Toate experimentele au rulat pe o singură maşină:

\begin{table}[h]
  \centering
  \caption{Mediul de antrenare şi evaluare.}
  \label{tab:hardware}
  \begin{tabular}{ll}
    \toprule
    GPU              & NVIDIA GeForce RTX 4090, 24\,GB GDDR6X, driver 590.48.01 \\
    CPU              & AMD EPYC-Milan, 8 vCPU (1 fir/nucleu) \\
    RAM              & 47\,GB; 24\,GB \texttt{tmpfs} folosit pentru pregătirea datelor \\
    Stocare          & volum rotaţional, $\approx$9\,MB/s secvenţial, $\approx$50 fişiere/s \\
    SO / Python      & Ubuntu (Linux 6.8), Python 3.10, Conda \\
    Framework        & PyTorch 2.4.0 (CUDA 11.8, cuDNN 9.1.0), torchvision 0.19.0 \\
    Biblioteci       & scikit-learn 1.7.2, NumPy 2.2.6, SciPy 1.15.2, Pillow 12.3.0 \\
    \bottomrule
  \end{tabular}
\end{table}

Înmulţirea de matrice/convoluţia TF32, autotunerul cuDNN şi precizia mixtă (AMP,
fp16) au fost activate pe tot parcursul. Etajul de entropie propriu-zis
\emph{nu} este eligibil pentru autocast --- \texttt{unfold}, \texttt{sort} şi
testele de egalitate nu au nuclee fp16 --- aşa că acesta se execută în fp32 şi
doar reţeaua convoluţională de bază rulează în semi-precizie.

\begin{quote}\small
\textbf{O notă despre I/O.} Volumul rădăcină al maşinii este rotaţional. Rulând
direct de pe disc, fiecare configuraţie era blocată la $\approx0.47$\,s/pas
indiferent de costul său real, deoarece GPU-ul stătea inactiv în timp ce
procesele de încărcare a datelor erau blocate în somn neîntreruptibil. Prin
urmare, toate duratele raportate provin din rulări al căror subset de date a
fost mai întâi transferat în \texttt{tmpfs}. Fără acest lucru, coloana de timp a
fiecărui baleiaj măsoară discul, nu modelul.
\end{quote}

\subsection{Datele}

Antrenarea a folosit corpusul standard ProGAN al lui Wang şi
colab.~\cite{wang2020cnn}, distribuit împreună cu CNNDetection/NPR: 20 de
categorii LSUN, fiecare cu un director \texttt{0\_real} şi unul
\texttt{1\_fake}. Dacă nu se precizează altfel, baleiajele au folosit opt
categorii (\texttt{car, cat, chair, horse, boat, person, dog, bird}), pentru a
evita un rezultat care să fie o afirmaţie despre maşini în particular. Evaluarea
a folosit partiţia \texttt{val} rezervată pentru baleiaje, iar pentru modelul
final întreaga suită de test ForenSynths / NPR: 13 generatoare GAN
(90{,}329 imagini) şi 13 generatoare de difuzie (51{,}992 imagini).

Imaginile sunt redimensionate la $256\times256$ şi decupate la $224\times224$ ---
\texttt{RandomCrop} pentru antrenare, \texttt{CenterCrop} pentru evaluare.

\begin{quote}\small
\textbf{Geometria de evaluare.} Fişierul \texttt{test.py} publicat (aici \texttt{mlep/upstream/test.py}) setează
\texttt{no\_crop=True}, alimentând cadre complete $256\times256$ la testare, deşi
antrenarea a folosit decupaje de $224$. Noi decupăm central la evaluare, astfel
încât statisticile BatchNorm reestimate pe decupaje de $224$
(Secţiunea~\ref{sec:setup-bn}) să fie aplicate la aceeaşi dimensiune a intrării.
Măsurat izolat, acest lucru valorează $\approx0.06$ AP --- mult mai puţin decât
efectul BatchNorm, dar gratuit de făcut corect.
\end{quote}

\subsection{Protocolul de antrenare}

Fiecare configuraţie dintr-un baleiaj se antrenează pentru exact acelaşi număr
de paşi de optimizare, astfel încât comparaţiile să fie corecte, iar timpul real
să difere doar în funcţie de costul pe pas al fiecărei configuraţii:

\begin{table}[h]
  \centering
  \caption{Protocolul de baleiaj. Rularea lungă finală
  (Secţiunea~\ref{sec:exp-degradation}) diferă şi este descrisă acolo.}
  \label{tab:protocol}
  \begin{tabular}{ll}
    \toprule
    Reţea de bază      & ResNet-18 (0.68M parametri) \\
    Paşi               & 3{,}000 (200 pentru testele de fum) \\
    Dimensiunea lotului& 16 (dimensionat pentru configuraţia cea mai costisitoare) \\
    Optimizator        & Adam, $\beta = (0.9, 0.999)$, rată de învăţare constantă $2\times10^{-4}$ \\
    Funcţia de cost    & \texttt{BCEWithLogitsLoss} pe un singur logit \\
    Decupaj            & $224$ dintr-o redimensionare la $256$, oglindire orizontală aleatoare \\
    Precizie           & AMP (reţea de bază fp16, etaj de intrare fp32) \\
    Recalibrare BN     & 50 de loturi de antrenare, medie cumulativă \\
    Validare           & 50 de loturi ($n=800$) per scenariu \\
    Seminţe            & 100, 200, 300 \\
    \bottomrule
  \end{tabular}
\end{table}

Dimensiunea lotului este menţinută \emph{identică} între configuraţii în mod
intenţionat. Argumentul de corectitudine (,,fiecare configuraţie se antrenează
acelaşi număr de paşi'') se susţine doar dacă un pas înseamnă acelaşi lucru
peste tot, aşa că lotul este dimensionat pentru configuraţia cea mai
costisitoare, în loc să fie reglat per configuraţie. Constrângerea determinantă
este memoria, nu viteza: calea generală \texttt{unfold}+\texttt{sort}
materializează un tensor de porţiuni $(B, 3, HW, w^2)$, iar \texttt{torch.sort}
adaugă indici int64 la dublul dimensiunii valorilor, astfel încât configuraţia
cu trei ferestre atinge un vârf de $21.3$\,GB din cei $24.6$\,GB ai plăcii la
lot 32.

\subsection{Scenariile de evaluare}

Fiecare model antrenat este evaluat în cinci condiţii de test, aplicate cu
probabilitatea $1.0$, astfel încât măsurătoarea să fie deterministă. Un generator
fix per încărcător de date garantează că fiecare scenariu vede aceleaşi imagini
de bază, deci diferenţele provin doar din corupţie:

\begin{center}
\begin{tabular}{ll}
  \toprule
  \texttt{clean}     & fără degradare \\
  \texttt{blur}      & gaussian, $\sigma = 2.0$ \\
  \texttt{jpeg}      & JPEG calitate 75 (PIL) \\
  \texttt{webp}      & WebP calitate 80 --- niciodată văzut la antrenare, o sondă în afara distribuţiei \\
  \texttt{blur+jpeg} & ambele, mai întâi blur \\
  \bottomrule
\end{tabular}
\end{center}

Raportăm precizia medie (AP) ca metrică principală, împreună cu ROC-AUC,
acurateţea la pragul $0.5$, cea mai bună acurateţe realizabilă la orice prag şi
logitul mediu. Ultimele două există pentru a diagnostica eşecul descris în
continuare.

\subsection{O precondiţie: recalibrarea BatchNorm}
\label{sec:setup-bn}

Primele noastre încercări de reproducere au produs un rezultat neinterpretabil:
acurateţea blocată la proporţia claselor ($\approx0.50$), în timp ce AP rămânea
ridicată, iar configuraţii diferite se prăbuşeau spre clase constante
\emph{diferite}. Cauza nu este modelul, ci contabilitatea BatchNorm.

Momentul implicit de $0.1$ al BatchNorm face din \texttt{running\_mean} şi
\texttt{running\_var} o medie mobilă exponenţială pe aproximativ ultimele 20 de
loturi --- circa 320 de imagini la lot 16 --- colectate în timp ce ponderile
încă se mişcau. Per strat, eroarea este mică (la \texttt{bn1} varianţa stocată
se situează cu circa 10\% sub varianţa reală a lotului), dar ea se acumulează
prin cele nouă straturi BatchNorm din \texttt{layer1}+\texttt{layer2}, iar capul
este un singur \texttt{Linear(128,1)} peste o agregare globală prin medie, deci
nimic în aval nu absoarbe această derivă.

Prin urmare, reestimăm toate statisticile BatchNorm după antrenare şi înainte de
evaluare, resetându-le, comutând momentul pe \texttt{None} (medie cumulativă) şi
rulând 50 de treceri înainte peste datele de \emph{antrenare}. Folosirea datelor
de validare aici ar fi transductivă. Efectul, la ponderi identice, este
dramatic:

\begin{table}[h]
  \centering
  \caption{Recalibrarea BatchNorm pe \texttt{baseline\_2x2}, 3{,}000 paşi, opt
  categorii, lot 16. Aceleaşi ponderi în ambele rânduri.}
  \label{tab:bn}
  \begin{tabular}{lccc}
    \toprule
     & logit mediu & acurateţe & AP \\
    \midrule
    \texttt{model.eval()} ca atare & $+24.16$ & 0.4938 & 0.8057 \\
    după recalibrare                & $-0.25$  & 0.9762 & 0.9985 \\
    \bottomrule
  \end{tabular}
\end{table}

Fiecare logit era deplasat cu aproximativ 24, prăbuşind fiecare predicţie într-o
singură clasă. Semnul este arbitrar --- o rulare de 600 de paşi s-a deplasat cu
$-13.85$ în schimb --- ceea ce explică exact de ce unele configuraţii raportau
totul ca real, iar altele totul ca fals. Esenţial, \textbf{şi AP se modifică}
($0.81 \to 0.9985$): aceasta nu a fost niciodată doar o problemă de prag, ci
corupea chiar ordonarea pe care baleiajele există să o producă. Fiecare cifră
raportată în acest capitol este de după recalibrare.

\section{Experimentul 1: dimensiunea ferestrei glisante}
\label{sec:exp-window}

\paragraph{Ipoteza.} MLEP fixează fereastra la $2\times2$, cea mai mică posibilă.
O fereastră $2\times2$ poate exprima doar cinci valori distincte de entropie
(Tabelul~\ref{tab:2x2}), ceea ce reprezintă o cuantizare foarte grosieră a
diversităţii locale. O fereastră mai mare admite un alfabet mai bogat ---
$w=4$ dă $K=16$ pixeli şi deci mult mai multe distribuţii distinctibile --- şi
ar putea capta amprente ale generatoarelor care operează pe o rază spaţială mai
mare.

\paragraph{Metoda.} Am antrenat câte o configuraţie pentru fiecare dimensiune de
fereastră $w \in \{2,3,4,5,6,7,8\}$, menţinând fixe piramida cu trei scări,
reţeaua de bază, numărul de paşi şi orice alt hiperparametru. Am testat
suplimentar \texttt{unique}, o măsură deliberat mai săracă decât entropia:
în loc să cântărească fiecare valoare după cât de des apare, ea numără pur şi
simplu câte valori distincte există în fereastră. Pe o fereastră $2\times2$
aceasta atribuie celor cinci cazuri din Tabelul~\ref{tab:2x2} valorile
$1, 2, 2, 3, 4$.

Acea repetiţie a lui $2$ nu este o scăpare, ci esenţa testului. Cazul ,,trei
egale'' şi cazul ,,două perechi'' conţin amândouă exact două valori distincte,
deci \texttt{unique} le \emph{contopeşte} într-o singură categorie, deşi
entropia Shannon le separă ($0.8113$ faţă de $1.0$). Cu alte cuvinte,
\texttt{unique} nu este o simplă reetichetare a celor cinci cazuri, ci o
cuantizare pe \emph{patru} niveluri în loc de cinci: distruge efectiv informaţie
pe care etajul de intrare o avea. Dacă distincţia dintre ,,trei egale'' şi
,,două perechi'' contează, \texttt{unique} trebuie să fie măsurabil mai slabă.

\begin{table}[t]
  \centering
  \caption{Baleiajul dimensiunii ferestrei. ResNet-18, 3{,}000 paşi, sămânţa 100,
  opt categorii. Precizie medie; mai mare este mai bine. ,,Timp'' este timpul
  real de antrenare în secunde.}
  \label{tab:window}
  \begin{tabular}{lrrrrrr}
    \toprule
    Configuraţie & clean & blur & jpeg & webp & blur+jpeg & timp \\
    \midrule
    \texttt{w2x2\_unique}  & 0.9980 & 0.6710 & 0.4144 & 0.4155 & 0.4218 & 42 \\
    \texttt{baseline\_2x2} & \textbf{0.9979} & 0.6817 & 0.4199 & 0.4193 & 0.4265 & 59 \\
    \texttt{w3x3}          & 0.9907 & 0.6548 & 0.4335 & 0.4232 & 0.4240 & 396 \\
    \texttt{w4x4}          & 0.9716 & 0.6094 & 0.4131 & 0.4064 & 0.4350 & 415 \\
    \texttt{w6x6}          & 0.8012 & 0.5654 & 0.4338 & 0.4236 & 0.4665 & 553 \\
    \texttt{w5x5}          & 0.7726 & 0.5663 & 0.4150 & 0.4119 & 0.4252 & 441 \\
    \texttt{w7x7}          & 0.5771 & 0.5116 & 0.4241 & 0.4259 & 0.4459 & 586 \\
    \texttt{w8x8}          & 0.5395 & 0.5195 & 0.4399 & 0.4402 & 0.4587 & 623 \\
    \bottomrule
  \end{tabular}
\end{table}

\paragraph{Rezultat: eşec.} AP pe date curate scade monoton odată cu dimensiunea
ferestrei, de la $0.998$ la $w=2$ până la $0.540$ la $w=8$ --- ultima abia peste
nivelul hazardului. Costul se mişcă în direcţia opusă: $w \geq 3$ părăseşte
tabela de căutare cu cinci intrări şi urmează calea generală bazată pe sortare,
făcând antrenarea cu un ordin de mărime mai lentă ($59$\,s faţă de $623$\,s)
pentru un model strict mai slab.

\paragraph{Interpretare.} Ipoteza a inversat mecanismul real. Amprenta
generatorului pe care o exploatează MLEP este o reţea de supraeşantionare cu o
perioadă de unul sau doi pixeli (Secţiunea~\ref{sec:worked-example}). O fereastră
$2\times2$ este \emph{potrivită} acestei perioade. Lărgirea ferestrei mediază
amprenta împreună cu textura autentică a scenei şi, deoarece entropia este o
statistică de mulţime, ea nu poate separa cele două contribuţii odată amestecate
--- o fereastră mare vede o distribuţie agitată indiferent dacă agitaţia provine
din reţeaua de supraeşantionare sau din iarbă. Cuantizarea grosieră pe cinci
niveluri nu este o limitare care trebuie depăşită; este un filtru adaptat.

\paragraph{Rezultatul testului \texttt{unique}.} Nu contează. \texttt{unique}
obţine $0.9980$ AP pe date curate faţă de $0.9979$ pentru Shannon --- o diferenţă
de $0.0001$, cu un ordin de mărime sub abaterea standard între seminţe de
$\pm0.002$ măsurată ulterior în Tabelul~\ref{tab:entropy-seeds}. Cele două sunt,
practic, aceeaşi configuraţie. Ba mai mult, \texttt{unique} este şi cu
$17$\,s mai rapidă, pentru că numărarea valorilor distincte este o operaţie mai
ieftină decât evaluarea unei entropii.

Acest rezultat este mai puternic decât pare la prima vedere. Nu am înlocuit
entropia cu o măsură echivalentă, ci cu una \emph{strict mai săracă}, care
contopeşte două dintre cele cinci cazuri --- şi nu s-a pierdut nimic. Prin
urmare, distincţia dintre ,,trei pixeli egali'' şi ,,două perechi'' este
irelevantă pentru sarcină: reţelei îi este suficient să ştie aproximativ câte
valori diferite se află într-o fereastră, iar modul exact în care acestea sunt
distribuite nu adaugă nimic.

Aceasta este prima indicaţie că \emph{funcţionala precisă} nu poartă informaţie
utilă dincolo de o partiţionare foarte grosieră a ferestrei --- şi ea prefigurează
direct rezultatul nul din Secţiunea~\ref{sec:exp-entropy}, unde şase funcţionale
suplimentare, alese tocmai pentru a repondera aceleaşi valori în moduri
diferite, ajung la aceeaşi concluzie. Merită observat că \texttt{unique} este de
fapt un membru al familiei Rényi testate acolo: la limita $\alpha \to 0$, entropia
Rényi devine exact $\log_2$ din numărul de valori distincte, adică
\texttt{unique} până la o transformare monotonă. Fără să fi intenţionat asta,
Experimentul 1 testase deja un capăt al intervalului pe care Experimentul 3 avea
să îl exploreze sistematic.

\section{Experimentul 2: ferestre multiple, normalizare şi înregistrare}
\label{sec:exp-multiwindow}

\paragraph{Ipoteza.} Deoarece nicio fereastră \emph{individuală} mai mare nu a
depăşit $2\times2$, întrebarea firească următoare este dacă mai multe ferestre
\emph{împreună} poartă informaţie complementară: harta $2\times2$ furnizează
amprenta, iar hărţile mai largi furnizează contextul care ar putea dezambiguiza
textura scenei de netezirea produsă de codec. Combinăm $w \in \{2,4,6\}$ peste
aceeaşi piramidă, obţinând $3 \times 3 \times 3 = 27$ canale de intrare.

Am identificat în plus două defecte ale combinaţiei naive şi am testat câte o
corecţie pentru fiecare, ca gemeni exacţi care schimbă un singur lucru.

\paragraph{Primul defect: nepotrivirea intervalului dinamic.} Valoarea maximă pe
care o poate lua entropia depinde de dimensiunea ferestrei. O fereastră
$2\times2$ conţine patru pixeli, deci cel mult patru valori distincte şi cel
mult $\log_2 4 = 2$ biţi, în timp ce o fereastră $6\times6$ conţine 36 de pixeli
şi urcă până la $\log_2 36 \approx 5.17$ biţi. Harta produsă de fereastra lată
are aşadar numere sistematic de circa $2.5$ ori mai mari --- nu pentru că ar fi
găsit mai multe dovezi despre generator, ci pur şi simplu pentru că fereastra
este mai mare. Când stivuim cele trei hărţi şi le dăm primei convoluţii, aceasta
înmulţeşte fiecare canal cu o pondere şi adună rezultatele, aşa că, la
iniţializare, canalul cu numere mari domină suma indiferent de cât de informativ
este. Iar acest dezechilibru nu se corectează mai târziu: \texttt{conv1} nu are
termen liber, iar \texttt{bn1} normalizează cele 64 de canale de \emph{ieşire},
adică după ce adunarea a avut deja loc, deci nu mai poate separa contribuţiile
individuale ale intrărilor. Corecţia este simplă: \texttt{normalize\_entropy}
împarte fiecare hartă la propriul ei maxim teoretic $\log_2 K$, unde $K = w^2$.
Toate hărţile ajung astfel în intervalul $[0,1]$ şi importanţa fiecăreia rămâne
să fie decisă de ponderile învăţate, nu de un accident al dimensiunii ferestrei.

\paragraph{Al doilea defect: dezalinierea spaţială.} O fereastră $w\times w$
glisată cu pasul $1$ încape de $H-w+1$ ori pe o latură, deci ferestrele mai mari
produc hărţi mai mici: la un decupaj de $224$ pixeli, fereastra $2\times2$ dă o
hartă de $223\times223$, cea $4\times4$ una de $221\times221$, iar cea $6\times6$
una de $219\times219$. Pentru a le putea stivui într-un singur tensor, ele
trebuie aduse la aceeaşi dimensiune, iar varianta implicită le reeşantionează
biliniar la dimensiunea celei mai mari. Problema este că reeşantionarea aliniază
marginile hărţilor şi le \emph{întinde} conţinutul, în loc să îl deplaseze:
valoarea de la poziţia $(i,j)$ descrie fereastra care \emph{începe} la $(i,j)$,
adică una centrată la $(i + (w-1)/2)$, aşa că hărţi cu ferestre diferite descriu
centre decalate între ele. Întinderea corectează acest decalaj în centrul
imaginii, dar îl lasă tot mai mare spre margini (până la $\approx1.6$ pixeli
pentru fereastra $6\times6$), iar interpolarea estompează în plus tocmai
structura de frecvenţă înaltă care constituie întregul semnal. O reţea
convoluţională poate absorbi o deplasare constantă în poziţiile filtrelor
învăţate, dar nu şi o deformare dependentă de poziţie, întrucât ponderile sunt
partajate.

Corecţia este să nu întindem nimic, ci să completăm hărţile mici cu
$d = (w-w_{\min})/2$ poziţii pe fiecare latură: harta $4\times4$ primeşte o
poziţie de fiecare parte ($221 + 2 = 223$), iar cea $6\times6$ două
($219 + 4 = 223$). Poziţia $(i,j)$ ajunge astfel să descrie exact acelaşi loc din
imagine în toate cele trei hărţi, fără nicio interpolare şi fără nicio pierdere
de detaliu. Valorile adăugate sunt copii ale marginii existente
(\texttt{mode='replicate'}), nu zerouri: pe o hartă de entropie, zero este o
valoare cu semnificaţie --- înseamnă ,,toţi pixelii ferestrei sunt identici'',
adică semnătura perfectă a unei zone supraeşantionate --- aşa că o bordură de
zerouri ar înconjura fiecare imagine cu un cadru artificial perfect plat, un
artefact puternic şi constant pe care reţeaua l-ar exploata imediat. Varianta
\texttt{pad} este exactă, fără interpolare, şi s-a măsurat că a costat
$+1$\,MiB şi niciun timp suplimentar.

Cele două corecţii sunt testate \emph{separat}: \texttt{multiwindow\_normalized}
schimbă doar normalizarea (păstrând reeşantionarea), iar
\texttt{multiwindow\_aligned} doar alinierea (păstrând scările neuniforme). Nicio
configuraţie nu le combină, aşa că un eventual efect al aplicării lor simultane
rămâne nemăsurat.

\begin{table}[t]
  \centering
  \caption{Configuraţii cu ferestre multiple faţă de modelul de referinţă.
  ResNet-18, 3{,}000 paşi, sămânţa 100.}
  \label{tab:multiwindow}
  \begin{tabular}{lrrrrrr}
    \toprule
    Configuraţie & clean & blur & jpeg & webp & blur+jpeg & timp \\
    \midrule
    \texttt{baseline\_2x2}                      & \textbf{0.9979} & 0.6817 & 0.4199 & 0.4193 & 0.4265 & 59 \\
    \texttt{multiwindow\_multiscale}            & 0.9965 & 0.6581 & 0.4189 & 0.4103 & 0.4210 & 962 \\
    \texttt{multiwindow\_normalized}            & 0.9938 & 0.6762 & 0.4153 & 0.4090 & 0.4242 & 963 \\
    \texttt{multiwindow\_aligned} (\texttt{pad})& 0.9944 & 0.6472 & 0.4146 & 0.4071 & 0.4194 & 1024 \\
    \bottomrule
  \end{tabular}
\end{table}

\paragraph{Rezultat: eşec.} Toate cele trei variante cu ferestre multiple sunt
marginal \emph{mai slabe} decât modelul de referinţă cu o singură fereastră
$2\times2$ pe date curate şi nu sunt mai bune sub nicio corupţie, la un cost de
antrenare de $16$--$17\times$ mai mare. Ambele corecţii principiale ---
normalizarea şi înregistrarea exactă --- au recuperat o parte din decalajul faţă
de combinaţia necorectată în unele celule (normalizarea îmbunătăţeşte blur de la
$0.6581$ la $0.6762$), dar niciuna nu a atins modelul de referinţă.

\paragraph{Interpretare.} Corecţiile erau reale; lucrul pe care îl corectau nu
merita avut. Adăugarea de canale cu ferestre late adaugă la \texttt{conv1}
intrări purtătoare de zgomot, iar nicio scalare şi înregistrare, oricât de
corectă, a unui canal neinformativ nu îl face informativ. Împreună cu
Experimentul 1, aceasta închide chestiunea \emph{geometriei} ferestrei: nici
dimensiunea ferestrei, nici numărul de dimensiuni de fereastră nu reprezintă
locul în care se află vreo îmbunătăţire.

\section{Experimentul 3: funcţionala de entropie}
\label{sec:exp-entropy}

\paragraph{Ipoteza.} Epuizând geometria, ne îndreptăm către funcţională.
Entropia Shannon ponderează contribuţia unei valori cu $-p\log_2 p$, ceea ce
este o alegere între multe altele. Entropiile generalizate reponderează valorile
rare faţă de cele frecvente printr-un exponent reglabil, iar o valoare rară este
exact ceea ce produce zgomotul unui senzor şi nu produce un supraeşantionator.
Concret, am testat:

\begin{align}
  \text{Rényi:}   \quad H_\alpha &= \frac{1}{1-\alpha}\log_2 \sum_g p_g^{\alpha},
  & \alpha &\in \{0.5, 2, 4\} \label{eq:renyi}\\
  \text{Tsallis:} \quad S_q      &= \frac{1 - \sum_g p_g^{q}}{q-1},
  & q &\in \{0.5, 2, 4\} \label{eq:tsallis}
\end{align}

\paragraph{Rényi: un buton care reglează cât contează valorile rare.} Entropia
Shannon dă fiecărei valori ponderea $-p\log_2 p$. Rényi introduce un exponent
reglabil $\alpha$ care mută atenţia măsurii: pentru $\alpha > 1$ domină valoarea
cea mai frecventă, iar măsura devine practic ,,cât de concentrată este
distribuţia''; pentru $\alpha < 1$ valorile rare capătă greutate, iar măsura se
apropie de ,,câte feluri de valori există, indiferent cât de des apar''; la
limita $\alpha \to 1$ se regăseşte exact Shannon. Un caz limită merită
menţionat: la $\alpha \to 0$, Rényi devine $\log_2$ din numărul de valori
distincte, adică exact măsura \texttt{unique} testată în
Secţiunea~\ref{sec:exp-window}. Rényi este aşadar o familie continuă care
conţine \texttt{unique} la un capăt şi Shannon la mijloc.

Motivul pentru care aceasta era o ipoteză rezonabilă ţine de mecanismul fizic:
zgomotul de senzor produce valori care apar o singură dată în fereastră, iar
supraeşantionarea nu produce aşa ceva. Accentuând valorile rare, ar trebui
amplificată tocmai diferenţa dintre real şi generat.

\paragraph{Tsallis: aceeaşi idee, fără logaritm.} Tsallis are şi ea un exponent
reglabil, $q$, şi porneşte de la aceeaşi sumă $\sum_g p_g^{q}$. Diferenţa este că
nu îi aplică logaritmul, ci o rescalare liniară. Consecinţa teoretică este că
Shannon şi Rényi sunt \emph{aditive} --- entropia a două sisteme independente
este suma entropiilor lor --- în timp ce Tsallis nu este, fiind construită în
fizica statistică exact pentru sisteme cu corelaţii pe distanţă lungă, unde
aditivitatea nu se mai susţine. Cum pixelii unei imagini sunt puternic coreloaţi,
alegerea părea potrivită. Consecinţa practică pe o fereastră de patru pixeli
este însă alta: se schimbă scara (maximul este $0.75$ în loc de $2$ biţi), dar
nu şi conţinutul informaţional.

\paragraph{De ce niciuna dintre ele nu putea ajuta.} Argumentul devine evident
numeric dacă evaluăm fiecare funcţională pe cele cinci cazuri posibile ale unei
ferestre $2\times2$ (Tabelul~\ref{tab:2x2}):

\begin{table}[h]
  \centering
  \caption{Valoarea atribuită de fiecare funcţională celor cinci configuraţii
  posibile ale unei ferestre $2\times2$. Toate cresc în aceeaşi ordine: nicio
  funcţională nu inversează vreo pereche şi niciuna nu separă cazuri pe care
  celelalte le confundă. Excepţia este \texttt{unique}, care \emph{contopeşte}
  cazurile $(3,1)$ şi $(2,2)$ --- şi care, cu toate acestea, obţine acelaşi scor
  (Secţiunea~\ref{sec:exp-window}).}
  \label{tab:functionals-2x2}
  \begin{tabular}{lcccccc}
    \toprule
    Caz & Shannon & Rényi $0.5$ & Rényi $2$ & Rényi $4$ & Tsallis $2$ & \texttt{unique} \\
    \midrule
    toate patru egale   & 0.000 & 0.000 & 0.000 & 0.000 & 0.000 & 1 \\
    trei egale $(3,1)$  & 0.811 & 0.900 & 0.678 & 0.548 & 0.375 & 2 \\
    două perechi $(2,2)$& 1.000 & 1.000 & 1.000 & 1.000 & 0.500 & 2 \\
    o pereche $(2,1,1)$ & 1.500 & 1.543 & 1.415 & 1.277 & 0.625 & 3 \\
    toate diferite      & 2.000 & 2.000 & 2.000 & 2.000 & 0.750 & 4 \\
    \bottomrule
  \end{tabular}
\end{table}

Toate coloanele cresc în aceeaşi ordine. Nicio funcţională nu inversează vreo
pereche de cazuri, niciuna nu uneşte două cazuri care erau separate şi niciuna
nu separă două cazuri care erau unite. Se schimbă doar \emph{spaţierea} celor
cinci niveluri --- iar prima convoluţie, urmată de BatchNorm, poate reaşeza
oricum acele cinci niveluri cum îi convine. Practic, reglăm din exterior un buton
pe care reţeaua îl are deja şi îl poate roti singură în timpul antrenării.

\paragraph{Permutare: singura cu adevărat diferită.} Entropia de permutare
(tipare ordinale Bandt--Pompe) nu se uită deloc la valori, ci doar la
\emph{ordinea} lor. Se iau trei pixeli consecutivi şi se reţine doar care este
cel mai mic, care este la mijloc şi care este cel mai mare; există $3! = 6$
tipare ordinale posibile. Se numără cât de des apare fiecare tipar şi se
calculează entropia acestei distribuţii de şase simboluri. Avantajul teoretic
este că măsura devine complet invariantă la orice transformare monotonă a
intensităţii --- luminozitate, gamma, contrast --- deci imună la o întreagă clasă
de manipulări fotometrice.

Aceasta este singura funcţională care aduce informaţie genuin nouă, întrucât
ordinea de rang nu este o funcţie a partiţiei de valori. Şi tot ea este cea care
pierde cel mai mult, dintr-un motiv care se vede imediat: \textbf{egalitatea este
invizibilă pentru o statistică de rang}. Dacă doi pixeli au exact aceeaşi
valoare, nu există ,,care este mai mare'' şi departajarea devine arbitrară --- iar
duplicarea valorilor \emph{este} întreaga amprentă exploatată de MLEP. Măsura
aruncă exact singurul indiciu care conta.

Întrucât o singură fereastră citită ca o secvenţă produce exact un tipar ordinal
(şi deci entropie zero), definim populaţia de tipare din interiorul unei ferestre
ca totalitatea subsecvenţelor de ordin 3 ale liniilor şi coloanelor sale; acest
lucru necesită $w \geq 3$. De aici decurge o rezervă pe care o semnalăm explicit:
entropia de permutare nu a putut fi niciodată evaluată la $w=2$, deci rezultatul
ei este parţial confundat cu efectul de dimensiune a ferestrei stabilit în
Secţiunea~\ref{sec:exp-window}, unde ferestrele mai mari pierdeau oricum.
Căderea măsurată este considerabil mai mare decât ar explica singur acel efect,
deci concluzia se menţine --- dar comparaţia nu este perfect controlată.

\paragraph{Metoda.} Experimentul s-a desfăşurat în două etape, astfel încât un
rezultat nul să nu poată fi pus pe seama optimizării modelului profund.

\begin{description}
  \item[Etapa de trăsături.] Statistici de rezumat ale entropiei (media,
    abaterea standard, minimul, maximul, mediana, IQR şi percentilele
    10/25/75/90 ale fiecărei hărţi locale), calculate pe imaginea întreagă şi pe
    ferestre neîntrepătrunse de $8\times8$/$16\times16$/$32\times32$, precum şi
    pe convenţia MLEP de $2\times2$, în tonuri de gri, R, G, B şi în vederea RGB
    comună. Patru clasificatori clasici --- regresie logistică, pădure aleatoare,
    SVM şi arbori cu gradient boosting --- au fost antrenaţi pe fiecare submulţime
    de ablaţie. Această etapă îşi poate permite ferestre pe imaginea întreagă şi
    de $32\times32$, care ar fi mult prea costisitoare ca etaj de intrare cu
    pasul 1.
  \item[Etapa profundă.] Reţeaua MLEP însăşi, antrenată cu fiecare etaj de
    entropie, prin exact aceeaşi funcţie \texttt{run\_config()} folosită de
    baleiajul de ferestre: aceiaşi paşi, aceeaşi recalibrare BN, aceleaşi
    scenarii de corupţie.
\end{description}

Grupurile de ablaţie au fost: A~Shannon (referinţă), B~doar Rényi, C~doar
Tsallis, D~doar permutare, E~Shannon+Rényi, F~Shannon+Tsallis,
G~Shannon+permutare, H~toate. Modurile combinate stivuiesc câte o hartă per
funcţională, canalele modului 0 fiind identice la nivel de bit cu configuraţia
corespunzătoare cu un singur mod, ceea ce face din fiecare combinaţie un test
\emph{aditiv} curat. În total, 46 de configuraţii au fost evaluate în rularea
combinată.

\begin{table}[t]
  \centering
  \caption{Etapa profundă, ROC-AUC, sămânţa 100. Sunt afişate doar rândurile de
  sus; tabelul complet cu 46 de rânduri apare în
  \texttt{results/entropy\_20260802\_160157\_seed100.txt}.}
  \label{tab:entropy-deep}
  \begin{tabular}{lrrrrr}
    \toprule
    Configuraţie & clean & blur & jpeg & webp & blur+jpeg \\
    \midrule
    \texttt{ent\_C\_tsallis4}  & 0.9987 & 0.7452 & 0.4365 & 0.4644 & 0.4644 \\
    \texttt{ent\_B\_renyi2}    & 0.9984 & 0.7408 & 0.4256 & 0.4316 & 0.4576 \\
    \texttt{ent\_C\_tsallis2}  & 0.9983 & 0.7435 & 0.4444 & 0.4432 & 0.4702 \\
    \texttt{ent\_B\_renyi4}    & 0.9982 & \textbf{0.7520} & 0.4253 & 0.4374 & 0.4595 \\
    \texttt{ent\_A\_shannon}   & 0.9978 & 0.7335 & 0.4282 & 0.4320 & 0.4574 \\
    \midrule
    \texttt{ent\_H\_all}       & 0.9854 & 0.7024 & 0.4317 & 0.4175 & 0.4688 \\
    \texttt{ent\_G\_shannon\_perm} & 0.9453 & 0.6481 & 0.4214 & 0.4002 & 0.4454 \\
    \texttt{ent\_D\_perm3}     & 0.8087 & 0.6038 & 0.4654 & 0.4390 & 0.4777 \\
    \texttt{ent\_D\_perm4}     & 0.7555 & 0.5913 & 0.4615 & 0.4576 & 0.4505 \\
    \bottomrule
  \end{tabular}
\end{table}

\paragraph{Rezultat: eşec (în limita zgomotului).} Ordonarea aparentă din
Tabelul~\ref{tab:entropy-deep} se întinde pe $0.9978$--$0.9987$ ROC-AUC pe date
curate, un interval de $0.0009$. Pentru a testa dacă ea rezistă, am reluat cele
cinci configuraţii fruntaşe pe trei seminţe:

\begin{table}[h]
  \centering
  \caption{Funcţionale de entropie, precizie medie, medie $\pm$ abatere standard
  pe seminţele 100/200/300.}
  \label{tab:entropy-seeds}
  \begin{tabular}{lccc}
    \toprule
    Configuraţie & clean & blur & jpeg \\
    \midrule
    \texttt{ent\_C\_tsallis4} & $0.997 \pm 0.002$ & $0.677 \pm 0.024$ & $0.448 \pm 0.023$ \\
    \texttt{ent\_B\_renyi4}   & $0.996 \pm 0.003$ & $0.684 \pm 0.015$ & $0.448 \pm 0.030$ \\
    \texttt{ent\_A\_shannon}  & $0.996 \pm 0.002$ & $0.669 \pm 0.014$ & $0.446 \pm 0.024$ \\
    \texttt{ent\_B\_renyi2}   & $0.996 \pm 0.002$ & $0.675 \pm 0.013$ & $0.446 \pm 0.031$ \\
    \texttt{ent\_C\_tsallis2} & $0.995 \pm 0.003$ & $0.673 \pm 0.015$ & $0.449 \pm 0.015$ \\
    \bottomrule
  \end{tabular}
\end{table}

Fiecare diferenţă este mai mică decât propria abatere standard. Singurul semnal
consistent este că Rényi cu $\alpha=4$ câştigă aproximativ $+0.015$ ROC-AUC pe
blur faţă de Shannon, în toate cele trei seminţe şi monoton în $\alpha$ ---
motiv pentru care l-am dus mai departe ca test explicit în
Secţiunea~\ref{sec:exp-degradation}, unde blurul este o ţintă, nu un efect
secundar. Nici acolo nu a rezistat.

Etapa de trăsături confirmă independent: cel mai bun clasificator per grup
depăşeşte Shannon cu cel mult $+0.0040$ ROC-AUC pe date curate (E, regresie
logistică), iar entropia de permutare singură \emph{pierde} $0.1074$.
Performanţa medie în condiţii corupte, pe toate grupurile, se situează în
$0.502$--$0.532$, adică abia peste nivelul hazardului pentru fiecare
funcţională.

\paragraph{Interpretare.} Acesta este cel mai puternic rezultat nul din capitol
şi decurge din Tabelul~\ref{tab:functionals-2x2}. La $w=2$ fereastra admite doar
cinci partiţii posibile ale celor patru valori. \emph{Orice} funcţională de
entropie este o funcţie doar de partiţie, deci fiecare funcţională din
Ecuaţiile~\ref{eq:renyi} şi \ref{eq:tsallis} este doar o reetichetare a
aceloraşi cinci categorii prin cinci numere reale diferite --- şi încă o
reetichetare strict monotonă. Prima convoluţie poate absorbi aproape complet o
rescalare monotonă a intrării sale. Schimbarea funcţionalei nu poate deci adăuga
informaţie; poate doar schimba spaţierea a cinci niveluri pe care reţeaua este
oricum liberă să le respaţieze. Entropia de permutare este singura funcţională
cu adevărat diferită, întrucât ordinea de rang nu este o funcţie de partiţia
valorilor --- şi este net \emph{mai slabă}, deoarece eliminarea magnitudinilor
elimină indiciul de duplicare a valorilor, care constituie întreaga amprentă.

Testul \texttt{unique} din Secţiunea~\ref{sec:exp-window} împinge concluzia cu
un pas mai departe. Rényi şi Tsallis nu puteau adăuga informaţie, fiindcă sunt
reetichetări monotone ale aceloraşi cinci cazuri; \texttt{unique}, în schimb,
chiar \emph{elimină} informaţie, contopind cazurile $(3,1)$ şi $(2,2)$ într-o
singură categorie --- şi nici asta nu a schimbat rezultatul. Aşadar, conţinutul
util al ferestrei este chiar mai grosier decât cele cinci niveluri: reţelei îi
este suficient să ştie aproximativ câte valori distincte există, iar modul în
care acestea sunt distribuite nu contează în niciun sens măsurabil.

(Acest argument al partiţiei produce şi un rezultat ingineresc util: o tabelă de
căutare cu cinci intrări, indexată după numărul de perechi egale, calculează
\emph{orice} entropie pe o fereastră $2\times2$, făcând Rényi şi Tsallis la fel
de ieftine ca Shannon --- $0.63$ faţă de $0.98$\,ms per hartă la lot 16, faţă de
$41$\,ms pentru calea generală.)

\section{Experimentul 4: entropia între canale (culoare comună)}
\label{sec:exp-color}

\paragraph{Ipoteza.} Experimentul 3 a stabilit că nicio reponderare scalară a
unei distribuţii de valori \emph{marginale} nu ajută. Dar MLEP îşi calculează
cele trei hărţi independent pe fiecare canal, deci etajul de intrare vede doar
trei marginale. O fereastră $2\times2$ poate fi maximal entropică în R, în G şi
în B şi totuşi să conţină doar două \emph{culori} distincte --- iar nicio
funcţie de marginale nu poate recupera acest lucru. Adăugăm prin urmare entropia
distribuţiei comune a tripletelor RGB, unde doi pixeli contează ca acelaşi
simbol doar dacă coincid pe toate cele trei canale.

Aceasta este motivată în mod specific de problema robusteţii. Atât JPEG cât şi
WebP convertesc în YCbCr, subeşantionează croma 4:2:0 şi cuantizează planurile
de cromă mult mai agresiv decât luma. Recompresia afectează deci structura
dintre canale şi structura per canal în măsuri \emph{diferite} --- o diferenţă
pe care un etaj de intrare bazat doar pe marginale nu are niciun canal în care
să o exprime.

\paragraph{Metoda.} Două variante: \texttt{joint} adaugă harta comună la
marginale ($4$ hărţi per (scară, fereastră) $=12$ canale, canalele 0--2 din
fiecare grup fiind identice la nivel de bit cu referinţa, deci un test strict
aditiv), iar \texttt{joint\_only} alimentează doar harta comună ($3$ canale,
\emph{mai puţine} decât referinţa), ca ablaţie care spune dacă statistica comună
poartă semnal în sine sau doar profită de context. Ambele au fost testate şi
peste antrenarea augmentată, întrucât acolo se află de fapt robusteţea.

\begin{table}[t]
  \centering
  \caption{Entropia comună de culoare. ResNet-18, 3{,}000 paşi, sămânţa 100.}
  \label{tab:color}
  \begin{tabular}{lrrrrr}
    \toprule
    Configuraţie & clean & blur & jpeg & webp & blur+jpeg \\
    \midrule
    \texttt{color\_joint\_2x2}          & \textbf{0.9982} & \textbf{0.6943} & 0.4063 & 0.4140 & 0.4247 \\
    \texttt{baseline\_2x2}              & 0.9979 & 0.6813 & 0.4194 & 0.4194 & 0.4269 \\
    \texttt{baseline\_train\_jpeg}      & 0.9866 & 0.6116 & 0.4726 & 0.4943 & 0.4510 \\
    \texttt{color\_joint\_train\_jpeg}  & 0.9834 & 0.6321 & \textbf{0.4759} & 0.4882 & \textbf{0.4911} \\
    \texttt{color\_jointonly\_2x2}      & 0.9821 & 0.5922 & 0.4919 & 0.4625 & 0.4730 \\
    \texttt{baseline\_train\_webp\_p1}  & 0.6462 & 0.5079 & 0.5910 & 0.5474 & 0.5246 \\
    \texttt{color\_joint\_train\_webp\_p1} & 0.4912 & 0.4524 & 0.5417 & 0.5552 & 0.5030 \\
    \bottomrule
  \end{tabular}
\end{table}

\paragraph{Rezultat: eşec.} Pe date curate harta comună valorează $+0.0003$ AP
($0.9982$ faţă de $0.9979$) --- imposibil de distins de dispersia între seminţe
de $\pm0.002$ măsurată în Tabelul~\ref{tab:entropy-seeds}. Sub corupţie este
\emph{mai slabă} decât referinţa pe jpeg ($0.4063$ faţă de $0.4194$). Combinată
cu antrenarea augmentată JPEG, produce într-adevăr cea mai mare valoare
\texttt{blur+jpeg} pe care am observat-o dintr-o schimbare a etajului de intrare
($0.4911$ faţă de $0.4510$), dar celula rămâne sub nivelul hazardului, deci
ordonarea este în continuare inversată --- modelul a devenit \emph{mai puţin
grav} inversat, nu corect. Ablaţia \texttt{joint\_only} pierde $0.016$ AP pe
date curate, arătând că statistica comună este mai slabă pe cont propriu decât
marginalele pe care le-ar înlocui, iar perechea antrenată pe WebP se prăbuşeşte
de-a dreptul la $0.4912$.

\paragraph{Interpretare.} Premisa era corectă --- statistica comună chiar nu
poate fi recuperată din marginale --- şi totuşi nu a ajutat. Duplicarea între
canale şi duplicarea per canal poartă, evident, informaţie aproape identică
pentru această sarcină, deoarece supraeşantionarea generatorului operează
simultan pe toate cele trei canale: un pixel duplicat este duplicat în R, G şi B
deodată, deci harta comună este, \emph{în practică}, aproape o funcţie
deterministă a marginalelor, chiar dacă în principiu nu este.

\section{Experimentul 5: separarea texturilor bogate/sărace}
\label{sec:exp-texture}

\paragraph{Ipoteza.} Fiecare schimbare a etajului de intrare de până acum a
eşuat, aşa că atacăm direct mecanismul de eşec identificat în
Secţiunea~\ref{sec:worked-example}. Problema este că o singură hartă de entropie
pe imaginea întreagă nu poate distinge ,,neted pentru că scena este netedă'' de
,,neted pentru că a netezit-o codecul''. Urmând PatchCraft~\cite{zhong2024patchcraft},
sortăm porţiunile imaginii după diversitatea texturii şi îi dăm reţelei regiunile
plate şi regiunile agitate ca două imagini \emph{separate}. Pânza săracă
izolează astfel exact ceea ce distruge un codec, pânza bogată ceea ce
supravieţuieşte în cea mai mare parte, iar reţeaua poate citi cele două populaţii
independent.

\paragraph{Metoda.} Diversitatea texturii unei porţiuni $p\times p$ este suma
fluctuaţiilor absolute ale pixelilor pe patru direcţii (orizontală, verticală,
diagonală, antidiagonală), calculată pe media RGB. Porţiunile sunt clasate,
împărţite în jumătatea superioară şi cea inferioară şi reasamblate în două pânze
de dimensiuni identice. Două moduri de combinare: \texttt{concat} stivuieşte
ambele teancuri de entropie (18 canale), iar \texttt{diff} alimentează diferenţa
lor (9 canale), ceea ce foloseşte PatchCraft însuşi şi este echivalent ca
capacitate cu referinţa, astfel încât o victorie acolo nu ar putea fi explicată
prin simplul fapt că \texttt{conv1} este mai lat.

Separarea rulează \emph{înainte} de amestecarea blocurilor şi pe imaginea
originală: diversitatea texturii are sens doar pe porţiuni coerente spaţial, iar
amestecarea prealabilă ar lăsa fiecare porţiune cu aceeaşi diversitate medie pe
imagine, reducând separarea la o departajare după ordinea de baleiere --- un
experiment care nu măsoară nimic, deşi arată perfect sănătos. Am ales $p=16$
deoarece grila de porţiuni are o axă pară la dimensiunile de decupaj 128, 224
\emph{şi} 256, aşa că separarea nu aruncă niciun pixel; valoarea $p=32$ folosită
de PatchCraft lasă o grilă $7\times7$ la decupaj 224 şi forţează o tăiere de o
coloană, care aruncă 14.3\% din decupaj.

\begin{table}[t]
  \centering
  \caption{Separarea texturilor. Blocul superior: antrenare pe date curate, 200
  de paşi, o categorie (triere). Blocul inferior: în pereche cu antrenarea
  augmentată, 3{,}000 de paşi, opt categorii.}
  \label{tab:texture}
  \begin{tabular}{lrrrrr}
    \toprule
    Configuraţie & clean & blur & jpeg & webp & blur+jpeg \\
    \midrule
    \texttt{baseline\_2x2}            & \textbf{0.9951} & 0.6717 & 0.5179 & 0.4866 & 0.5047 \\
    \texttt{texsplit\_2x2}            & 0.9646 & 0.6208 & 0.4888 & 0.4710 & 0.4863 \\
    \texttt{texsplit\_2x2\_diff}      & 0.9550 & 0.6497 & 0.4978 & 0.4865 & 0.5085 \\
    \midrule
    \texttt{texsplit\_2x2\_train\_blur} & \textbf{0.9931} & 0.6879 & 0.4188 & 0.4174 & 0.4225 \\
    \texttt{baseline\_train\_blur}      & 0.9926 & \textbf{0.7209} & 0.4177 & 0.4152 & 0.4175 \\
    \texttt{baseline\_train\_jpeg}      & 0.9865 & 0.6119 & \textbf{0.4726} & 0.4940 & 0.4508 \\
    \texttt{texsplit\_2x2\_train\_jpeg} & 0.9855 & 0.6092 & 0.4616 & 0.4903 & 0.4410 \\
    \bottomrule
  \end{tabular}
\end{table}

\paragraph{Rezultat: eşec.} Antrenată pe date curate, separarea costă $0.030$ AP
pe date curate (\texttt{concat}) sau $0.040$ (\texttt{diff}) şi nu ajută în
nicio celulă coruptă. În pereche cu antrenarea augmentată --- celula în care ar
fi trebuit să se vadă, întrucât acolo se află robusteţea --- valorează $+0.0005$
pe blur şi $-0.0010$ pe jpeg, ambele mult sub pragul de zgomot, \emph{pierzând}
totodată $0.033$ AP pe blur faţă de \texttt{baseline\_train\_blur}.

\paragraph{Interpretare.} Atribuim acest lucru unei interacţiuni cu amestecarea
blocurilor, specifică MLEP. Separarea din PatchCraft funcţionează pe o reţea care
consumă direct porţiuni; MLEP amestecă în plus blocuri $2\times2$, ceea ce
distruge deja coerenţa spaţială care face distincţia bogat/sărac semnificativă în
aval. Cele două operaţii sunt amestecări parţial redundante, iar costul separării
--- înjumătăţirea suportului spaţial efectiv al fiecărei pânze --- este plătit
fără beneficiul ei. Configuraţia de control dedicată
(\texttt{baseline\_2x2\_norearr}, care dezactivează amestecarea) a fost
implementată pentru a izola acest lucru, dar un baleiaj complet în pereche cu ea
rămâne nerealizat şi constituie singurul capăt liber al acestui experiment.

\section{Experimentul 6: antrenare cu augmentare prin degradări}
\label{sec:exp-aug}

\paragraph{Ipoteza.} Întrucât nicio schimbare a \emph{reprezentării} nu a
restabilit robusteţea, testăm remediul standard: arătarea de imagini degradate
modelului în timpul antrenării. Dacă etajul de entropie păstrează vreun semnal
utilizabil sub compresie, augmentarea ar trebui să permită clasificatorului să
îl găsească.

\paragraph{Metoda.} Două augmentări, fiecare aplicată cu probabilitatea $0.5$ în
timpul antrenării: blur gaussian cu $\sigma \sim U(0,3)$ şi JPEG la calitatea 30
sau 100, folosind fie codificatorul OpenCV, fie pe cel PIL. Evaluarea rămâne
neschimbată.

O primă rundă a arătat fiecare celulă AP coruptă \emph{sub} $0.5$ (0.415--0.494,
12/12 celule pe patru configuraţii), în timp ce acurateţea se situa la
$\approx0.52$. După cum s-a stabilit în Secţiunea~\ref{sec:worked-example}, un AP
sub nivelul hazardului înseamnă o ordonare inversată. Am emis ipoteza că
probabilitatea de $0.5$ era ea însăşi cauza: aceasta antrenează pe un amestec
50/50 de imagini curate şi corupte ale căror ordonări de clasă sunt
\emph{opuse}, iar domeniul curat este separabil în proporţie de $\approx0.99$
faţă de $\approx0.5$ pentru cel corupt, aşa că domeniul curat domină gradientul
şi ordonarea sa este aplicată invers pe intrarea coruptă. Am reluat prin urmare
rularea cu probabilitatea $1.0$, antrenând pe un singur domeniu în loc de un
amestec contradictoriu.

\begin{table}[t]
  \centering
  \caption{Antrenare augmentată. ResNet-18, 3{,}000 paşi, sămânţa 100.
  \texttt{\_p1} = probabilitate de corupţie 1.0 în loc de 0.5.}
  \label{tab:aug}
  \begin{tabular}{lrrrrr}
    \toprule
    Configuraţie & clean & blur & jpeg & webp & blur+jpeg \\
    \midrule
    \texttt{baseline\_2x2}             & \textbf{0.9979} & 0.6813 & 0.4191 & 0.4191 & 0.4267 \\
    \texttt{baseline\_train\_blur}     & 0.9936 & \textbf{0.7229} & 0.4148 & 0.4153 & 0.4170 \\
    \texttt{baseline\_train\_jpeg}     & 0.9866 & 0.6092 & 0.4726 & 0.4923 & 0.4533 \\
    \texttt{baseline\_train\_jpeg\_p1} & 0.9339 & 0.4927 & 0.4747 & 0.4557 & 0.4295 \\
    \texttt{baseline\_train\_webp\_p1} & 0.6431 & 0.5079 & \textbf{0.5950} & \textbf{0.5484} & \textbf{0.5297} \\
    \bottomrule
  \end{tabular}
\end{table}

\paragraph{Rezultat: parţial şi revelator.} Augmentarea cu blur este singurul
succes neechivoc din acest capitol: ridică AP pe blur de la $0.6813$ la $0.7229$
la un cost de doar $0.004$ AP pe date curate. Dar câştigul este strict limitat la
degradarea pe care s-a antrenat --- jpeg şi webp rămân neschimbate sau uşor mai
slabe.

Augmentarea JPEG ridică AP pe jpeg de la $0.4199$ la $0.4726$ şi pe webp de la
$0.4193$ la $0.4923$ şi este singura configuraţie al cărei logit mediu rămâne
aproape de zero în fiecare scenariu; dar ambele celule rămân \emph{sub} $0.5$.
Modelul este mai puţin grav inversat, nu corect.

Testul cu probabilitate 1.0 a confirmat ipoteza amestecului şi, simultan, a
demolat-o ca soluţie. \texttt{baseline\_train\_webp\_p1} este singura
configuraţie din întregul studiu care împinge \emph{vreo} celulă coruptă
decisiv peste nivelul hazardului (jpeg $0.5950$, webp $0.5484$, blur+jpeg
$0.5297$) --- şi o face distrugând performanţa pe date curate, care scade de la
$0.998$ la $0.6431$. Nu există niciun punct de funcţionare intermediar: modelul
poate detecta imagini generate \emph{sau} poate supravieţui recompresiei, dar
etajul de entropie nu admite o reprezentare care să le facă pe amândouă.

\paragraph{Interpretare.} Aceasta este cea mai clară formulare a problemei
centrale. Amprenta imaginii curate şi amprenta imaginii recomprimate nu sunt doar
diferite, ci \emph{antagonice}: o reprezentare optimizată pentru una este
anti-optimizată pentru cealaltă. Exact acest lucru îl prezice exemplul lucrat din
Secţiunea~\ref{sec:worked-example}, iar nicio ponderare a augmentării nu le
reconciliază, deoarece conflictul este în trăsătură, nu în distribuţia de
antrenare.

\section{Experimentul 7: poate MLEP măcar să raporteze degradarea?}
\label{sec:exp-degradation}

\paragraph{Ipoteza.} Dat fiind că MLEP nu poate detecta \emph{prin} degradare,
punem o întrebare diferită: poate acelaşi etaj de intrare să \emph{identifice}
degradarea? Un răspuns pozitiv ar fi util într-un flux de prelucrare --- un
detector care ştie că priveşte o imagine puternic recomprimată se poate abţine
sau poate transfera decizia unui al doilea model --- şi măsoară totodată direct
câtă informaţie specifică degradării poartă efectiv hărţile de entropie.

\paragraph{Metoda.} Detectorul AI-vs-real nu este rescris. Păstrăm etajul de
intrare propriu al MLEP la setările sale ($2\times2$ ferestre, piramidă cu trei
scări) şi doar lărgim stratul final de la un logit la trei:
\[
  z_0 = \Pr(\text{generat de IA}), \quad
  z_1 = \Pr(\text{neclarizat}), \quad
  z_2 = \Pr(\text{comprimat JPEG}),
\]
antrenaţi împreună cu o singură funcţie \texttt{BCEWithLogitsLoss}. Întrucât
\texttt{num\_classes} este deja un argument al ResNet, \emph{nu} există nicio
schimbare arhitecturală.

Degradarea este aplicată pe o grilă deterministă $7 \times 4$ în loc să fie
eşantionată: calitatea JPEG $\in \{\text{fără pierderi}, 15, 30, 45, 60, 75, 100\}$
şi blurul gaussian $\in \{$fără $0.0$, slab $0.5$, mediu $1.5$, puternic $3.0\}$,
blurul înaintea JPEG. Calitatea 100 se consideră comprimată; doar imaginea
neatinsă este ,,necomprimată''. Patru celule --- $(30,\text{mediu})$,
$(75,\text{slab})$, $(15,\text{puternic})$, $(100,\text{mediu})$ --- sunt
excluse complet din antrenare şi formează condiţia \texttt{unseen}, un test de
\emph{combinaţie} nevăzută.

Celula pe care o primeşte un eşantion este o funcţie deterministă de calea sa
sursă, $\text{crc32}(\text{cale}) + 7919v + \text{sămânţă} \bmod |\text{celule}|$,
în loc de o extragere aleatoare. Acest lucru face atribuirea reproductibilă şi,
esenţial, \emph{independentă de eticheta de clasă}, astfel încât modelul nu poate
folosi degradarea ca scurtătură pentru ,,generat''. Fiecare variantă a unei surse
rămâne în interiorul partiţiei acelei surse; ambele proprietăţi sunt verificate
prin teste unitare, iar fiecare rulare înregistrează o verificare explicită de
scurgere a datelor (206{,}752 surse de antrenare, niciuna prezentă în vreuna
dintre cele şase mulţimi de evaluare).

\paragraph{Corecţia echilibrului claselor.} Lista simplă de antrenare cu 24 de
celule conţine exact o celulă neatinsă, deci o sursă ajunge curată în
$1/24 \approx 4\%$ din cazuri, iar capetele de blur/JPEG se antrenează pe
71\%/83\% exemple pozitive. Acest lucru a înfometat capul AI. Am supraeşantionat
celula curată de $12\times$, obţinând 36 de intrări dintre care 13 sunt
neatinse. Aceasta corectează bine marginalele:

\begin{center}
\begin{tabular}{lcc}
  \toprule
   & înainte & după \\
  \midrule
  fracţiunea curată      & 4.2\%  & 36.1\% \\
  pozitive blur          & 70.8\% & 47.2\% \\
  pozitive JPEG          & 83.3\% & 55.6\% \\
  \bottomrule
\end{tabular}
\end{center}

\paragraph{Rularea finală de antrenare.} Spre deosebire de baleiaje, această
configuraţie a fost antrenată până la convergenţă:

\begin{table}[h]
  \centering
  \caption{Modelul final conştient de degradare.}
  \label{tab:final-training}
  \begin{tabular}{ll}
    \toprule
    Reţea de bază      & varianta ResNet-50 a etajului de intrare, 1.44M parametri, 9 canale de intrare \\
    Iniţializare       & pornire ,,caldă'' din \texttt{pretrained/model\_epoch\_best.pth}, \texttt{fc1} reiniţializat \\
    Paşi / lot         & 200{,}000 $\times$ 32 (AMP), $\approx$7.7 epoci \\
    Optimizator        & Adam $(0.9, 0.999)$, fără regularizare a ponderilor \\
    Program lr         & lr $5\times10^{-5}$ cosinusoidal după 2{,}000 de paşi de încălzire, prag inferior $10^{-6}$ \\
    Date               & 206{,}752 surse $\times$ 4 variante $\approx$ 827k eşantioane, 8 categorii \\
    Timp real          & 9{,}886\,s ($\approx$2\,h\,45\,m) pe RTX 4090 \\
    Post-antrenare     & BatchNorm recalibrat pe 50 de loturi, \emph{apoi} salvat \\
    \bottomrule
  \end{tabular}
\end{table}

\begin{table}[t]
  \centering
  \caption{Experimentul de degradare, ROC-AUC per condiţie. Rândurile cu 3{,}000
  de paşi sunt rulările din protocolul de baleiaj; ultimul rând este modelul cu
  200{,}000 de paşi.}
  \label{tab:degradation}
  \begin{tabular}{llrrrrrr}
    \toprule
    Model & Cap & clean & doar jpeg & doar blur & ambele & seen & unseen \\
    \midrule
    3k paşi, Shannon     & AI   & 0.8750 & 0.5602 & 0.6679 & 0.5746 & 0.5743 & 0.5411 \\
    3k paşi, Rényi-4     & AI   & 0.9209 & 0.5928 & 0.6915 & 0.5780 & 0.5907 & 0.5389 \\
    200k paşi, Shannon   & AI   & \textbf{0.9999} & 0.7141 & 0.8898 & 0.6401 & 0.7281 & 0.6099 \\
    200k paşi, Shannon   & blur & --- & --- & --- & --- & 0.8705 & --- \\
    200k paşi, Shannon   & JPEG & --- & --- & --- & --- & \textbf{1.0000} & --- \\
    \bottomrule
  \end{tabular}
\end{table}

\paragraph{Rezultat: sarcina auxiliară reuşeşte, cea principală nu.} Capul JPEG
atinge un ROC-AUC de $1.0000$ cu o matrice de confuzie perfectă, iar capul de
blur $0.8705$. Capul AI atinge $0.9999$ pe imagini curate --- confirmând că
remedierea echilibrului claselor a funcţionat --- dar se degradează la $0.7141$
sub JPEG şi la $0.6401$ când ambele degradări sunt prezente.

Pe suita completă de test rezervată, modelul final obţine, pe fiecare mulţime de
generatoare:

\begin{center}
\begin{tabular}{lrrr}
  \toprule
  Mulţime & generatoare & acurateţe medie & AP mediu \\
  \midrule
  Mulţimea de difuzie & 13 & 82.8\% & 92.4\% \\
  Mulţimea GAN 1      & 13 & 72.7\% & 80.8\% \\
  \bottomrule
\end{tabular}
\end{center}

cu scoruri aproape perfecte pe generatoarele apropiate de datele de antrenare
(ProGAN 99.3\% acurateţe / 100.0 AP, StarGAN 99.9/100.0, StyleGAN2 94.4/100.0) şi
eşecuri pe altele (CRN 48.3/35.4, SeeingDark 50.0/40.6, improved-diffusion
55.6/49.6). Un baleiaj dedicat de perturbaţii confirmă că decalajul de robusteţe
supravieţuieşte antrenării complete: StarGAN scade de la AP 100.0 pe date curate
la 58.7 la calitate JPEG 75, o cădere de 41.3 puncte --- pe o degradare care era
\emph{în} grila de antrenare.

\paragraph{Interpretare şi rezerve.} Două concluzii. În primul rând, hărţile de
entropie chiar \emph{poartă} informaţie aproape completă despre faptul că o
imagine a fost comprimată JPEG --- deci incapacitatea de a detecta prin compresie
nu este o lipsă de informaţie despre compresie, ci o incapacitate de a o
factoriza în afara deciziei AI-vs-real. În al doilea rând, şi negativ,
antrenarea multi-sarcină nu a adus niciun câştig de robusteţe pentru capul
principal.

Menţionăm cinstit că scorul perfect al capului JPEG este parţial un artefact al
protocolului nostru: blurul rulează întotdeauna \emph{înaintea} JPEG, deci
compresia este întotdeauna ultima operaţie asupra pixelilor, iar artefactele sale
de blocuri $8\times8$ ajung la reţea intacte, niciodată estompate de o
reeşantionare ulterioară. O ordine aleatorizată a operaţiilor ar face sarcina
substanţial mai grea. Sonda WebP este informativă aici: capul JPEG se activează
la $\approx1.0$ pe WebP, care nu a fost niciodată în antrenare, deci ceea ce a
învăţat este ,,este comprimată'', nu ,,este JPEG''.

O a doua limitare priveşte condiţia \texttt{unseen}. Toate cele patru celule
excluse au \emph{atât} JPEG $\neq$ fără pierderi, \emph{cât şi} blur $\neq$ fără,
deci acea condiţie nu conţine exemple negative pentru niciunul dintre capetele
auxiliare, iar AP/ROC-AUC ale acestora sunt acolo nedefinite. Mulţimea exclusă ar
fi trebuit să includă cel puţin o celulă cu o singură degradare şi una curată;
aşa cum a fost proiectată, \texttt{unseen} măsoară doar capul AI.

\section{Sinteza experimentelor}
\label{sec:summary}

\begin{table}[t]
  \centering
  \caption{Toate modificările investigate. $\Delta_{\text{clean}}$ este
  variaţia AP pe date curate, iar $\Delta_{\text{corr}}$ variaţia AP mediată pe
  cele patru scenarii corupte, ambele relative la controlul potrivit sub acelaşi
  protocol. Rândurile marcate cu $^\dagger$ sunt împerecheate cu referinţa
  \emph{augmentată} corespunzătoare, nu cu \texttt{baseline\_2x2}; rândurile
  marcate cu $^\ddagger$ provin din rularea de triere de 200 de paşi. Abaterea
  standard între seminţe este $\pm0.002$ pe date curate şi până la $\pm0.03$ pe
  date corupte, deci orice se află în această bandă este zgomot.}
  \label{tab:summary}
  \small
  \begin{tabular}{llrrl}
    \toprule
    \# & Modificare & $\Delta_{\text{clean}}$ & $\Delta_{\text{corr}}$ & Rezultat \\
    \midrule
    1 & Fereastră $3\times3$                     & $-0.0072$ & $-0.0030$ & eşec \\
    1 & Fereastră $8\times8$                     & $-0.4584$ & $-0.0223$ & eşec (monoton în $w$) \\
    1 & \texttt{unique} în loc de Shannon        & $+0.0001$ & $-0.0062$ & niciun efect (deşi pierde un nivel) \\
    \midrule
    2 & Ferestre multiple $\{2,4,6\}$            & $-0.0014$ & $-0.0098$ & eşec \\
    2 & \quad + normalizare per fereastră        & $-0.0041$ & $-0.0057$ & eşec \\
    2 & \quad + înregistrare exactă (pad)        & $-0.0035$ & $-0.0148$ & eşec \\
    \midrule
    3 & Rényi $\alpha \in \{0.5,2,4\}$           & $\pm0.001$ & $+0.015$ (blur) & în limita zgomotului \\
    3 & Tsallis $q \in \{0.5,2,4\}$              & $\pm0.001$ & $\pm0.010$ & în limita zgomotului \\
    3 & Entropie de permutare$^{*}$              & $-0.1891$ & $+0.037$ & eşec \\
    3 & Toate funcţionalele stivuite$^{*}$       & $-0.0124$ & $\pm0.010$ & eşec \\
    \midrule
    4 & Entropie comună de culoare (aditivă)     & $+0.0003$ & $-0.0020$ & în limita zgomotului \\
    4 & Doar entropie comună de culoare          & $-0.0158$ & $+0.0180$ & eşec \\
    4 & Entropie comună + augm. JPEG$^\dagger$   & $-0.0032$ & $+0.0144$ & în limita zgomotului \\
    \midrule
    5 & Separare texturi, concat$^\ddagger$      & $-0.0309$ & $-0.0277$ & eşec \\
    5 & Separare texturi, diff$^\ddagger$        & $-0.0403$ & $-0.0090$ & eşec \\
    5 & Separare texturi + augm. blur$^\dagger$  & $+0.0005$ & $-0.0061$ & în limita zgomotului \\
    5 & Separare texturi + augm. JPEG$^\dagger$  & $-0.0010$ & $-0.0068$ & în limita zgomotului \\
    \midrule
    6 & Augmentare blur ($p=0.5$)                & $-0.0043$ & $\mathbf{+0.0056}$ & \textbf{succes, îngust} \\
    6 & Augmentare JPEG ($p=0.5$)                & $-0.0113$ & $+0.0196$ & parţial, celulele tot $<0.5$ \\
    6 & Augmentare JPEG ($p=1.0$)                & $-0.0640$ & $-0.0237$ & eşec \\
    6 & Augmentare WebP ($p=1.0$)                & $-0.3548$ & $\mathbf{+0.0584}$ & compromis, nu o soluţie \\
    \midrule
    7 & Predicţie a degradării cu mai multe capete & $\approx 0$ & $\approx 0$ & doar sarcină auxiliară \\
    \bottomrule
  \end{tabular}
  \\[2pt]
  \raggedright\footnotesize
  $^{*}$ ROC-AUC în loc de AP (etapa profundă a entropiei raportează ROC-AUC).
\end{table}

Pentru succesele înguste merită izolată celula specifică, nu media. Augmentarea
cu blur mută scenariul \texttt{blur} de la $0.6817$ la $0.7229$ ($+0.0412$);
augmentarea JPEG mută \texttt{jpeg} de la $0.4199$ la $0.4726$ ($+0.0527$);
augmentarea WebP la $p=1.0$ mută \texttt{jpeg} de la $0.4199$ la $0.5950$
($+0.1751$), singura valoare coruptă decisiv peste nivelul hazardului din studiu,
cu preţul a $0.355$ AP pe date curate.

\paragraph{Reproductibilitate.} \texttt{baseline\_2x2} a fost reantrenat de la
zero în trei rulări independente, în zile diferite, şi a returnat un AP pe date
curate de $0.9979$ în toate cele trei, cu un AP corupt mediu de $0.4869$,
$0.4866$ şi $0.4868$. Aparatul de măsură este deci stabil la aproximativ
$\pm0.0003$ pentru o sămânţă fixată, ceea ce justifică tratarea diferenţelor de
sub $0.002$ de mai sus ca zgomot, nu ca efecte mici.

\subsection{Ameninţări la adresa validităţii}

Trei rezerve limitează forţa acestor rezultate negative.

\begin{enumerate}
  \item \textbf{Durata antrenării.} Baleiajele folosesc 3{,}000 de paşi ca
    aproximare pentru antrenarea completă. O modificare care are nevoie de mai
    mult timp pentru a da roade ar fi punctată aici ca eşec. Singura
    configuraţie pe care am antrenat-o 200{,}000 de paşi
    (Secţiunea~\ref{sec:exp-degradation}) s-a îmbunătăţit substanţial în termeni
    absoluţi --- ROC-AUC pe date curate $0.875 \to 0.9999$ --- dar decalajul său
    \emph{relativ} de robusteţe nu s-a închis, ceea ce constituie o dovadă slabă
    că ordonarea este stabilă în raport cu durata antrenării.
  \item \textbf{Seminţele.} Doar comparaţia entropiilor a fost rulată pe trei
    seminţe. Diferenţele cu o singură sămânţă sub $\approx0.002$ AP pe date
    curate nu sunt interpretabile şi ne-am abţinut de la a revendica vreuna.
  \item \textbf{Domeniul grilei de augmentare.} Blurul este întotdeauna aplicat
    înaintea JPEG, iar reîncodarea după redimensionare --- fluxul cel mai
    frecvent în lumea reală --- nu este testată.
\end{enumerate}

\section{Concluzie şi tranziţie}
\label{sec:conclusion}

Am investigat unsprezece modificări distincte ale etajului de entropie al MLEP,
acoperind geometria ferestrei (dimensiune, multiplicitate, normalizare,
înregistrare), funcţionala de entropie (Rényi, Tsallis, permutare şi
combinaţiile lor), structura canalelor (entropia comună de culoare între
canale), descompunerea spaţială (separarea texturilor bogate/sărace) şi
distribuţia de antrenare (augmentare cu blur, JPEG şi WebP la două
probabilităţi). Nu am reuşit să găsim o modificare care să îmbunătăţească
acurateţea de detecţie a MLEP pe imagini curate dincolo de zgomotul dintre
seminţe şi nu am reuşit să găsim niciuna care să îi restabilească robusteţea la
recompresie.

Rezultatele converg către o singură explicaţie, ilustrată concret în
Secţiunea~\ref{sec:worked-example}. Trăsătura MLEP este o numărare a duplicării
locale de valori. Supraeşantionarea generativă creează o astfel de duplicare;
la fel şi compresia cu pierderi. Cele două cauze nu sunt doar confundabile ---
ele acţionează asupra trăsăturii în \emph{aceeaşi direcţie}, astfel încât o
fotografie reală puternic comprimată prezintă aceeaşi semnătură de entropie
scăzută ca o imagine generată. Acesta este motivul pentru care fiecare precizie
medie în condiţie coruptă pe care am măsurat-o se află sub $0.5$: nu un detector
slab, ci unul inversat. Şi este motivul pentru care fiecare intervenţie a eşuat
în acelaşi mod. Schimbarea ferestrei schimbă scara spaţială la care se măsoară
ambiguitatea; schimbarea funcţionalei reetichetează aceleaşi cinci partiţii;
adăugarea hărţii comune de culoare adaugă o statistică pe care
supraeşantionarea o corelează oricum cu marginalele. Niciuna dintre ele nu
adaugă un canal în care ,,neted pentru că este generat'' şi ,,neted pentru că
este comprimat'' să fie \emph{separabile}, deoarece această distincţie nu este
deloc prezentă în duplicarea locală de valori. Experimentul 7 exprimă acest
lucru cu precizie: acelaşi etaj de intrare identifică compresia JPEG cu un
ROC-AUC de $1.0000$, dar nu poate folosi acea cunoaştere pentru a-şi proteja
propria decizie AI-vs-real.

Singura intervenţie care a mutat decisiv o celulă coruptă peste nivelul
hazardului (augmentarea WebP la probabilitatea $1.0$, AP pe jpeg $0.5950$) a
făcut-o sacrificând AP pe date curate de la $0.998$ la $0.643$. Acest compromis
este cea mai clară dovadă că limitarea este de natură reprezentaţională, nu una
de antrenare: în interiorul acestui spaţiu de trăsături, detecţia pe imagini
curate şi detecţia robustă la compresie sunt obiective antagonice.

Concluzionăm, prin urmare, că robusteţea MLEP nu poate fi reparată din interiorul
propriului său etaj de intrare şi că informaţia lipsă trebuie să provină dintr-o
reprezentare \emph{diferită} --- una ale cărei moduri de eşec nu sunt aliniate cu
cele ale MLEP. Acest lucru motivează abordarea combinată investigată în restul
lucrării, în care MLEP este împerecheat cu un al doilea detector construit pe
principii diferite, astfel încât degradările care distrug dovezile unui model să
le lase intacte pe ale celuilalt. Această combinaţie este dezvoltată în capitolul
următor.
